\documentclass[8pt,a4paper,landscape]{extarticle}
\usepackage[UTF8]{ctex}
\usepackage[margin=0.35cm]{geometry}
\usepackage{multicol}
\usepackage{amsmath,amssymb,mathtools}
\usepackage{enumitem}
\usepackage{titlesec}
\usepackage{array,booktabs}
\usepackage{xcolor}
\usepackage{hyperref}

\setlength{\columnsep}{0.18cm}
\setlength{\columnseprule}{0.1pt}
\setlength{\parindent}{0pt}
\setlength{\parskip}{0pt}
\setlist{nosep,leftmargin=*,topsep=0pt,partopsep=0pt,parsep=0pt,itemsep=0pt}
\titlespacing*{\section}{0pt}{1.2pt}{0.4pt}
\titlespacing*{\subsection}{0pt}{0.8pt}{0.2pt}
\titleformat{\section}{\bfseries\scriptsize\color{blue!55!black}}{}{0pt}{}
\titleformat{\subsection}{\bfseries\tiny\color{black}}{}{0pt}{}
\pagenumbering{gobble}
\renewcommand{\arraystretch}{0.82}
\setlength{\tabcolsep}{2pt}
\raggedcolumns

\newcommand{\key}[1]{\textbf{#1}}
\newcommand{\term}[1]{\textit{#1}}
\newcommand{\R}{\mathbb{R}}
\newcommand{\E}{\mathbb{E}}
\newcommand{\N}{\mathcal{N}}
\newcommand{\dd}{\,\mathrm{d}}
\DeclareMathOperator*{\argmax}{arg\,max}
\DeclareMathOperator*{\argmin}{arg\,min}
\DeclareMathOperator{\softmax}{softmax}
\DeclareMathOperator{\IoU}{IoU}

\begin{document}
\tiny
\begin{multicols*}{4}

\section{L1 Introduction}
\key{Computer Vision (CV)}: 用计算机理解 image/video 的 meaning, 即 Marr: know what is where by looking. 图像来自 sensing device, CV 是 graphics 的 inverse problem: graphics $\text{model}\to\text{image}$, vision $\text{image}\to\text{world/model}$. \key{Understand image}: objects/scene, 3D/material properties, actions, relations, 可回答关于图像的问题。
\key{Applications}: measurement (shape/motion capture, 3D urban modeling, camera ruler, photo tourism), recognition (face/OCR/fine-grained/action), reconstruction, robotics, AR/VR, medical, autonomous driving. \key{Core difficulties}: inverse problem is ill-posed; projection loses depth; appearance changes with illumination/viewpoint/occlusion/deformation/background; semantic gap from pixels to meaning.
\key{Task modes}: measurement (shape/motion/3D), recognition (semantic content), manipulation/generation (inpainting, style/domain transfer, deepfake). Progress is dataset-driven: larger/richer data enables captioning, interaction, egocentric/video and multimodal tasks.

\section{L2 Image Basics \& Filters}
\subsection{Digital Image}
Digital camera: sensor array converts photons to electrons. Image formation includes sampling continuous image plane and quantizing intensity. Grayscale image $I:\R^2\to\R$; digital image $I\in\R^{H\times W\times C}$, usually uint8 $[0,255]$. Resolution is spatial sampling density. Low sampling causes aliasing; quantization causes intensity error.
\key{Image as function}: transformation can be pointwise/intensity, e.g. $g(x,y)=f(x,y)+50$, geometric $g(x,y)=f(x,-y)$, threshold/mask $g=f\cdot\mathbf 1[f<250]$.
\key{Color}: perceived property from light/surface/visual system. RGB is linear primary weights; HSV separates hue/saturation/value, useful when color appearance and brightness need separate handling. Histogram counts intensities/color bins; useful for contrast, thresholding, matching, normalization.
\key{Noise}: salt-and-pepper = random black/white spikes; impulse = sparse white spikes; Gaussian = iid normal intensity perturbation. Mean/Gaussian smoothing assumes neighboring pixels are similar and noise is independent, so averaging suppresses noise but blurs edges.

\subsection{Correlation, Convolution, Filtering}
\key{Correlation}: $G=H\otimes F$, $G[i,j]=\sum_{u,v}H[u,v]F[i+u,j+v]$. \key{Convolution}: kernel flipped, $(H*F)[i,j]=\sum_{u,v}H[u,v]F[i-u,j-v]$. For symmetric kernels, correlation = convolution. Boundary: zero fill, replicate border, mirror reflect, mirror across edge.
\key{Linear filters}: shifting, blur, sharpen, edge-like high-pass. Box filter averages but creates blocky response. Gaussian
$G_\sigma(x,y)=\frac{1}{2\pi\sigma^2}\exp(-(x^2+y^2)/(2\sigma^2))$; smoothing suppresses high frequency/noise; larger $\sigma$ means stronger blur. Median filter is nonlinear, robust to salt-and-pepper noise.
\key{Separable filter}: $K=uv^\top$, so 2D convolution becomes horizontal then vertical; cost $O(k^2HW)\to O(2kHW)$. Gaussian is separable. Sharpening can be viewed as $I+\alpha(I-\mathrm{blur}(I))$.
\key{Kernel facts}: smoothing kernels usually positive and sum to $1$, preserving constant images. Output shape choices: full keeps all overlaps; same keeps input size; valid only positions where kernel fully fits. \key{Convolution properties}: shift-invariant, commutative, associative, distributive, identity impulse $\delta$, and $\mathcal F(f*g)=\mathcal F(f)\mathcal F(g)$; Gaussian blur is low-pass. \key{Template matching}: filters can be matched templates; use normalized cross-correlation to reduce brightness bias. \key{Hybrid image}: low-pass one image + high-pass another, perceived content changes with viewing distance/scale. Separable iff $\operatorname{rank}(K)=1$; SVD gives best rank-1 approximation for nearly separable filters.

\section{L3 Camera Models}
\subsection{Perspective Projection}
Camera center $C$, optical axis $Z$, image plane distance $f$, principal point $p=(p_x,p_y)$. For camera coordinate $Q=(X,Y,Z)^\top$:
\[
x=\frac{fX}{Z}+p_x,\qquad y=\frac{fY}{Z}+p_y .
\]
Depth ambiguity: one pixel corresponds to a whole ray. Larger focal length $\Rightarrow$ narrower field of view.
\key{Pinhole/lens}: tiny aperture gives sharp image but little light and diffraction; large aperture gives more light but blur; lens focuses rays from one depth plane onto sensor. Projection loses depth, metric distances and angles; it preserves incidence of points/lines, but objects through camera center may degenerate.
\key{Homogeneous coordinates}: $(x,y,1)^\top\sim(sx,sy,s)^\top$. Line $l=(a,b,c)^\top$ means $l^\top x=0$; line through two points $l=x_1\times x_2$. Intrinsics:
\[
K=\begin{bmatrix}f_x&s&p_x\\0&f_y&p_y\\0&0&1\end{bmatrix},\qquad
\tilde x\sim KX_c .
\]
Square/no skew: $s=0,f_x=f_y=f$. Full projection:
\[
X_c=RX_w+t=R(X_w-C),\quad P=K[R\mid t],\quad \tilde x\sim PX_w .
\]
After $\tilde x=PX$, divide by third coordinate.

\subsection{Projective Properties}
Points map to points; 3D lines generally map to image lines; parallel 3D lines may meet at vanishing point. For line $X(t)=V+tD$, as $t\to\infty$, projected point depends only on direction $D$:
\[
\tilde x_v\sim KD,\qquad D\sim K^{-1}\tilde x_v .
\]
Vanishing points of directions in a plane lie on the \term{vanishing line}; ground-plane vanishing line is horizon. Parallel planes share horizon line. Orthographic projection drops depth:
\[
(X,Y,Z,1)^\top\mapsto(X,Y,1)^\top
\]
with no perspective division; good when depth variation is small relative to distance.
\key{VP caveats}: lines parallel to image plane meet at infinity; 2D image intersection does not imply 3D parallel lines. \key{Calibration}: known 3D calibration pattern + detected 2D corners estimates $K$ and extrinsics; three orthogonal vanishing points can calibrate Manhattan scenes. \key{Single-view metrology}: with flat ground and level camera, horizon height can reveal camera height on vertical objects; with known $K$ and ground plane, backproject image rays to plane and place CAD/measure scene.

\section{L4 Stereo \& Epipolar Geometry}
\subsection{Parallel Calibrated Stereo}
Two views turn a left pixel ray into an epipolar line in right image; matching + triangulation recovers 3D. For rectified parallel cameras with baseline $B$, focal $f$:
\[
y_l=y_r,\qquad d=x_l-x_r,\qquad Z=\frac{fB}{d}.
\]
Closer objects have larger disparity. Full 3D from left image:
\[
X=\frac{(x_l-p_x)Z}{f},\qquad Y=\frac{(y_l-p_y)Z}{f}.
\]
Patch matching cost for disparity $d$:
\[
\mathrm{SSD}(p,d)=\sum_{q\in W}(I_l(p+q)-I_r(p-(d,0)+q))^2.
\]
Normalized correlation robust to intensity shifts:
\[
\mathrm{NC}=\frac{\sum (I_l-\bar I_l)(I_r-\bar I_r)}
{\sqrt{\sum(I_l-\bar I_l)^2}\sqrt{\sum(I_r-\bar I_r)^2}}.
\]
Patch size: small = detailed/noisy; large = stable/blurred boundaries. Smooth stereo energy:
\[
E(d)=\sum_pD_p(d_p)+\lambda\sum_{(p,q)\in N}V(d_p,d_q).
\]

\subsection{General Epipolar Geometry}
Epipole $e_r$: left camera center projected into right image. Epipolar plane $(C_l,C_r,P)$ intersects images in paired epipolar lines; matching search becomes 1D. Fundamental matrix maps point to line:
\[
l_r=Fp_l,\qquad p_r^\top Fp_l=0.
\]
Construction: $F=[e_r]_\times H_\pi$, where $H_\pi$ is plane-induced homography and
\[
[e]_\times=\begin{bmatrix}0&-e_3&e_2\\e_3&0&-e_1\\-e_2&e_1&0\end{bmatrix}.
\]
Each correspondence gives linear equation in $F$ entries:
\[
\begin{aligned}
&x_rx_lf_{11}+x_ry_lf_{12}+x_rf_{13}
+y_rx_lf_{21}+y_ry_lf_{22}\\
&\quad +y_rf_{23}+x_lf_{31}+y_lf_{32}+f_{33}=0 .
\end{aligned}
\]
$F$ has scale ambiguity and rank constraint; use multiple matches + least squares/RANSAC. Rectification warps images so epipolar lines become horizontal. From $F$ one projective camera choice: $P_l=[I\mid0]$, $P_r=[[e_r]_\times F\mid e_r]$.
\key{F facts}: $F$ has scale ambiguity + rank-2 constraint, so 7 DoF; $e_r^\top F=0$ and $Fe_l=0$. Need $n\ge7$ matches in principle, more with RANSAC in practice. Learned stereo can replace SSD/NCC with CNN/classifier cost. \key{SfM}: multi-view stereo jointly estimates cameras and 3D points by minimizing reprojection error $\sum_{ij}\|x_{ij}-\pi(P_iX_j)\|^2$; solved by bundle adjustment. \key{Structured light}: projector emits known pattern to simplify correspondence; active stereo can use one camera + projector, e.g. Kinect IR.

\section{L5 Keypoints, Features, Matching}
\key{Local feature goal}: match corresponding image points from same world point. Recall needs invariance to rotation/scale/affine/illumination; precision needs contrast/distinctiveness and locality. Raw patch descriptor is sensitive to shifts/rotations.
\key{Detector desiderata}: salient + repeatable under viewpoint/illumination. Matching all pixels increases false positives unless descriptor is extremely strong.
\subsection{SIFT \& Keypoint Selection}
\key{SIFT}: Scale-Invariant Feature Transform. Scale invariance from Gaussian pyramid/scale-space; illumination robustness from normalized gradients; rotation invariance from dominant orientation. Descriptor: rotate local grid to dominant gradient, split into $4\times4$ cells, each with 8-bin gradient histogram $\Rightarrow128$D descriptor.
\key{Lowe ratio test}: match nearest neighbor in descriptor space only if $r=d_1/d_2$ is low; high ratio means best and second-best matches are similarly plausible, so the match is ambiguous. Modern alternatives: supervised matchers learn from paired images (e.g. SuperGlue-style); SSL/VLM/diffusion features (DINO/CLIP/SigLIP) provide semantic correspondences.
\key{Corners}: intensity changes in all directions. Harris uses local second-moment matrix:
\[
\Delta^2\approx
\begin{bmatrix}\Delta x&\Delta y\end{bmatrix}
\begin{bmatrix}D_x^2&D_xD_y\\D_xD_y&D_y^2\end{bmatrix}
\begin{bmatrix}\Delta x\\\Delta y\end{bmatrix}.
\]
Large eigenvalues in both directions $\Rightarrow$ corner; one large one small $\Rightarrow$ edge; both small $\Rightarrow$ flat. Harris response often $R=\det M-k(\operatorname{tr}M)^2$. Scale-invariant interest points: Laplacian of Gaussian approximated by Difference of Gaussian.

\subsection{Homography \& RANSAC}
Transformations: translation/rotation/similarity/affine/projective. Affine preserves parallel lines; projective/homography may send parallel lines to vanishing point. Planar scene or pure camera rotation gives homography:
\[
\begin{bmatrix}x'\\y'\\1\end{bmatrix}\sim
H\begin{bmatrix}x\\y\\1\end{bmatrix},\quad
H=\begin{bmatrix}a&b&c\\d&e&f\\g&h&1\end{bmatrix}.
\]
$H$ has 8 DoF up to scale; each match gives 2 equations; $\ge4$ matches, over-constrained solve least squares. Image mosaic: reproject images from same center of projection into a shared synthetic wide-angle camera.
\key{2D transform hierarchy}: translation 2 DoF; similarity 4 DoF (translation+rotation+uniform scale); affine 6 DoF, preserves parallel lines; homography/projective 8 DoF, preserves straight lines but not parallelism. \key{Warping}: forward warp can create holes/non-integer hits; inverse sampling with interpolation is usually cleaner. Mosaic pipeline: detect/match features $\to$ estimate $H$ with RANSAC $\to$ warp to common frame $\to$ blend.
\key{RANSAC}: sample minimal set, fit model, count inliers under threshold, repeat, refit on inliers. Iteration count for success probability $p$, inlier ratio $w$, sample size $s$:
\[
N=\frac{\log(1-p)}{\log(1-w^s)}.
\]
Pros: simple, works with outliers; cons: poor with low inlier ratio, bad thresholds, degenerate samples.

\section{L6 Vision Neural Networks}
\subsection{CNN Basics}
CNN biases: local connectivity, weight sharing, spatial structure. Conv filter $F\times F\times C$ slides over input; output depth = number of filters $K$.
\[
\begin{gathered}
W_2=\frac{W_1-F+2P}{S}+1,\quad
H_2=\frac{H_1-F+2P}{S}+1,\\
\#\theta=F^2CK+K.
\end{gathered}
\]
Pooling: per-channel downsample, no parameters; max pooling $W_2=(W_1-F)/S+1$. $1\times1$ conv mixes channels only, useful for projection/bottleneck. ReLU $f(x)=\max(0,x)$: efficient, non-saturating positive side, but dead ReLU/output not zero-centered.
\key{FC vs Conv}: FC flattens $HWC$ into a vector; params for $M$ outputs are $HWC\cdot M+M$, no spatial sharing. Conv keeps the grid, uses local receptive fields and shared weights. Sigmoid/tanh saturate, kill gradients, are not zero-centered (sigmoid), and use expensive exp; ReLU is fast/non-saturating for $x>0$ but zero-gradient for $x<0$; LeakyReLU/ELU mitigate dead units.
\key{Initialization/normalization}: constant init breaks symmetry poorly; small random can vanish/explode. Xavier std $1/\sqrt{D_{in}}$. LayerNorm:
\[
\mu=\frac1d\sum_jx_j,\quad \sigma^2=\frac1d\sum_j(x_j-\mu)^2,\quad y_j=\gamma\frac{x_j-\mu}{\sqrt{\sigma^2+\epsilon}}+\beta.
\]
Regularization/generalization: early stopping, data augmentation, dropout, stochastic depth, random crop, Mixup. \key{Regularization view}: always keep best validation snapshot; augmentation/dropout/stochastic depth can be read as training with injected noise, while test-time prediction marginalizes/averages over the noise.

\subsection{Architectures}
AlexNet: ReLU, data augmentation, GPU split, ImageNet 2012. ResNet solves deep optimization by residual block $H(x)=F(x)+x$; if $F=0$ identity is easy. Bottleneck: $1\times1$ reduce, $3\times3$, $1\times1$ expand. DenseNet: each layer connects to all later layers, feature reuse, gradient flow. Depthwise separable conv decomposes standard conv into depthwise spatial + pointwise channel:
\[
K^2C^2\to K^2C+C^2.
\]
ConvNeXt: Transformer ideas in CNN: patchify stem, depthwise conv, inverted bottleneck, GeLU, LayerNorm, separate downsampling. U-Net: encoder-decoder with long skip connections for localization; descendants in segmentation, diffusion, point cloud. \key{ResNet head}: drops large FC stack; last conv $\to$ global average pooling $\to$ class FC, reducing params and supporting variable spatial input. \key{ConvNeXt V2}: not just architecture; ConvNeXt blocks are co-designed with MAE-style/self-supervised training.

\subsection{Attention, Transformer, ViT}
Attention with query $Q$, data $X$: $K=XW_K,V=XW_V$,
\[
E=\frac{QK^\top}{\sqrt{D}},\quad A=\softmax(E),\quad Y=AV.
\]
Self-attention: $Q=XW_Q,K=XW_K,V=XW_V$. Masked attention sets illegal future logits to $-\infty$. Multi-head: $H$ heads, $D_H=D/H$, concatenate then output projection. Complexity/memory $O(N^2)$ due to attention matrix; FlashAttention avoids materializing full $N\times N$ matrix.
Transformer block: self-attention is vector interaction; LayerNorm/MLP are per-token; pre-norm makes residual identity easier. ViT: split image into patches, flatten, linear project to tokens, add positional embedding, Transformer, pool/[CLS], classify. Example $16\times16\times3=768$ patch vector. ViT weakens CNN translation equivariance but gives global interaction.
RMSNorm:
\[
y_i=\frac{x_i}{\sqrt{\epsilon+\frac1N\sum_i x_i^2}}\gamma_i.
\]

\section{L7 Classification \& Representation Learning}
\subsection{Recognition \& Self-Supervision}
Recognition tasks: identification, scene/object classification, annotation, detection, segmentation, pose, attributes, action. Classification challenge: semantic gap, viewpoint, clutter, illumination, occlusion, deformation, intra-class variation. Supervised pipeline: collect labels, train classifier, evaluate; expensive labels motivate self-supervised learning (SSL).
\key{Pretext vs downstream}: labels generated from data itself for pretraining; downstream is real labeled task. \key{Pretext risk}: auto-generated labels may teach task artifacts rather than transferable semantics if the pretext task is too narrow.
\key{MAE}: mask high ratio patches (e.g. 75\%), encoder sees visible patches only, decoder reconstructs pixels at masked patches; keep encoder, discard decoder. Loss:
\[
\mathcal L_{\mathrm{MAE}}=\frac1{|\mathcal M|}\sum_{i\in\mathcal M}\|x_i-\hat x_i\|_2^2.
\]
Asymmetric: large encoder, small decoder.

\subsection{Contrastive Learning}
Goal: positive pairs close, negatives far. InfoNCE:
\[
L=-\E\log\frac{\exp(s(f(x),f(x^+)))}{\exp(s(f(x),f(x^+)))+\sum_{j=1}^{N-1}\exp(s(f(x),f(x_j^-)))}.
\]
SimCLR: two augmentations of same image form positive; cosine score $s(u,v)=u^\top v/(\|u\|\|v\|)$; encoder $h=f(\tilde x)$, projection head $z=g(h)$. NT-Xent for positive $(i,j)$:
\[
\ell_{ij}=-\log\frac{\exp(s(z_i,z_j)/\tau)}{\sum_{k\ne i}\exp(s(z_i,z_k)/\tau)}.
\]
Downstream uses $h$ not $z$. MoCo: queue of keys decouples negatives from batch size; key encoder EMA:
\[
\theta_k\leftarrow m\theta_k+(1-m)\theta_q.
\]
DINO: student matches teacher outputs; teacher EMA, centering+sharpening avoid collapse:
\[
\theta_t\leftarrow \tau\theta_t+(1-\tau)\theta_s.
\]
DINOv2 scales curated data + global DINO + local masked/iBOT, strong patch features. \key{SSL eval}: pretrain encoder, discard projection head, then freeze+linear probe or fine-tune on downstream labels. \key{Energy view}: contrastive methods lower energy for positives and raise energy for negatives; regularized methods restrict low-energy volume / latent variance to avoid collapse without explicit negatives.

\subsection{JEPA \& Vision-Language}
JEA aligns embeddings $s_x=Enc_x(x),s_y=Enc_y(y)$. JEPA predicts target representation:
\[
\hat s_y=Pred(Enc_x(x),z),\qquad \mathcal L=\|\hat s_y-s_y\|^2.
\]
Not pixel reconstruction; predicts abstract representation, avoids irrelevant texture. Since $Enc_y$ is not invertible, JEPA does not define normalized $p(y|x)$ and is not a standard generative/probabilistic model. H-JEPA passes latents hierarchically to predict more abstract/longer-term structure. Need information constraints: $Enc_x,Enc_y$ informative; prediction error small; latent $z$ low-information so it cannot store answer. Collapse: all inputs map to constant; SIGReg forces latent distribution variance/normality.
\key{CLIP}: image-text contrastive pretraining on pairs, open vocabulary. Batch similarity:
\[
I_i=f_{\mathrm{img}}(x_i),\quad T_j=f_{\mathrm{text}}(t_j),\quad S_{ij}=I_i^\top T_j.
\]
Diagonal matched, off-diagonal negatives; train image-to-text and text-to-image. Zero-shot:
\[
\hat c=\argmax_{c\in\mathcal C} f_{\mathrm{img}}(x)^\top f_{\mathrm{text}}(\text{``a photo of a }c\text{''}).
\]
SigLIP replaces softmax contrastive with sigmoid pairwise loss, less batch-size coupled. CoCa combines CLIP-style contrastive with captioning decoder. \key{CLIP limitations}: image/text embeddings form separated cones (modality gap); image-level captions are weak on relations, counting, and compositional order (horse eating grass vs grass eating horse), improved by hard negatives, region captions, or stronger VLM training. LLaVA connects visual encoder to LLM via projection and instruction tuning; Flamingo uses gated cross-attention, masked attention, and in-context multimodal learning.

\section{L8 Object Detection}
\subsection{Task \& Metrics}
Detection output set $\{(b_i,c_i,s_i)\}_{i=1}^N$: boxes, class, confidence. Harder than classification: variable count/scale/location, overlaps, background. Approaches: sliding window, region proposal, implicit shape/codebook, modern learned queries. \key{Implicit Shape Model}: local features match codebook entries; each match casts vote for object center/scale, Hough-style detection and rough segmentation.
\key{IoU}: $\IoU(A,B)=\frac{|A\cap B|}{|A\cup B|}$. Precision/recall:
\[
P=\frac{TP}{TP+FP},\quad R=\frac{TP}{TP+FN}.
\]
AP = area under PR curve after sorting detections by confidence; mAP = mean over classes/IoU thresholds. Each GT matched once; duplicate boxes become FP. Anchor/proposal box regression often uses
\[
t_x=\frac{x-x_a}{w_a},\quad t_y=\frac{y-y_a}{h_a},\quad t_w=\log\frac{w}{w_a},\quad t_h=\log\frac{h}{h_a}.
\]
Detection loss:
\[
\mathcal L=\mathcal L_{cls}+\lambda\mathcal L_{box}.
\]

\subsection{R-CNN Family}
R-CNN: Selective Search $\sim2000$ proposals, crop/warp each, CNN per proposal, classify + box correction; slow due to repeated convolution. Fast R-CNN: one whole-image backbone, crop shared feature with RoIPool, heads. RoIPool maps variable RoI to fixed $7\times7$ by quantized bins; quantization hurts alignment. Faster R-CNN: learn proposals with RPN. For feature map $H\times W$ and $K$ anchors/location: candidates $KHW$, objectness logits $2KHW$, box offsets $4KHW$; sort by objectness and keep top $\sim300$ proposals. Two-stage = RPN proposals then RoI classify/regress.
\key{Single-stage}: YOLO/SSD/RetinaNet dense predict on grid. YOLO output:
\[
S\times S\times(5B+C),
\]
where each box has objectness + $(x,y,w,h)$ and classes. NMS: keep highest score, remove lower boxes with high IoU, repeat.
Final YOLO class-specific score is roughly $P(\mathrm{obj})P(c)$; cell/grid assignment makes dense detection fast but can struggle with crowded small objects.

\subsection{DETR \& Open Vocabulary}
DETR: CNN/ViT features + positional encoding + Transformer encoder + learned object queries in decoder; each query predicts class (incl. no-object) and box. Set prediction uses Hungarian matching:
\[
\hat\sigma=\argmin_\sigma\sum_{j=1}^M \mathcal C(y_j,\hat y_{\sigma(j)}),
\]
cost = class + L1 box + GIoU. Matching encourages one prediction per GT, reducing NMS need. Tradeoff: clean end-to-end but slow convergence/small objects.
Open-vocabulary detection replaces fixed class head by text queries. OWL-ViT: CLIP-style pretraining then detection fine-tuning; removes CLIP pooling so per-patch/token image embeddings compare to text/image queries, while a box MLP predicts coordinates. Text embedding acts as classifier weights. OWLv2/self-training expands detection supervision. Grounding DINO: language-aware open-set detector; bridges text phrase to boxes.

\section{L9 Dense Prediction}
\subsection{Segmentation, FCN, Upsampling}
Dense prediction: $I\in\R^{H\times W\times3}\to Y\in\R^{H\times W\times C}$. Semantic segmentation = class per pixel; instance = distinguish object instances; panoptic = stuff + things. Depth $D\in\R^{H\times W}$; flow $C=2$; diffusion predicts dense noise/clean image. Core tension: semantic abstraction vs spatial precision.
FCN converts classifier to fully convolutional: $1\times1$ conv over feature map gives low-res class score map; bilinear or learned transposed-conv upsampling restores resolution. Skip fusion combines high-level semantic coarse maps with low-level fine maps; FCN fuses skip logits by summing. Segmentation metrics:
\[
IoU_c=\frac{TP_c}{TP_c+FP_c+FN_c},\qquad mIoU=\frac1C\sum_cIoU_c.
\]
Upsampling: interpolation, transposed convolution (linear operator transpose, not inverse), max unpooling with switches, resize+conv to reduce checkerboard. Stride-2 transposed conv can be viewed as inserting zero rows/cols then convolving with flipped kernel; uneven overlap causes checkerboard. DeconvNet: encoder-decoder; U-Net concatenates encoder features with decoder features and predicts at the end; FPN: top-down + lateral connections to semantic pyramid.

\subsection{Mask R-CNN, SAM, Depth, Generation}
Mask R-CNN = Faster R-CNN + FCN mask head on RoIs. Outputs $\{(b_i,c_i,m_i)\}$. Mask head $M\in[0,1]^{K\times m\times m}$, uses per-pixel sigmoid BCE for GT class. RoIAlign avoids RoIPool quantization using bilinear sampling; crucial for masks/keypoints. Keypoints as heatmaps $\hat H\in\R^{K\times H\times W}$.
SAM: promptable segmentation foundation model. Components: heavy ViT image encoder (cacheable), prompt encoder (points/boxes/text/masks), lightweight mask decoder with two-way attention. Outputs multiple masks for ambiguity + predicted IoU for ranking. Limitations: may miss fine structures, hallucinate small disconnected components, fail boundaries/text-to-mask, and lose to domain-specific tools. SA-1B built with model-in-the-loop stages; SAM 2 stores previous-frame memory for video, preserving object identity under motion/occlusion and supporting streaming interaction.
Depth Anything: teacher/student, labeled + 65M unlabeled pseudo labels, DINOv2 backbone; perturbation consistency (color jitter/blur/CutMix) lets student generalize beyond teacher. V2 uses synthetic data for precise depth then bridges sim-to-real. Single-image depth is not true 3D and not necessarily multi-view consistent; DA3 adds depth+ray+pose/cross-view geometry targets.
Generation as dense prediction: GAN/cGAN/Pix2Pix/DDPM/LDM/DiT all predict dense maps over pixels/latents/patches. PatchGAN outputs spatial realism map over local patches; LDM denoises latent $z_t$; DiT replaces U-Net denoiser with ViT.

\section{L10 Synthetic Data}
\key{Why}: controllable labels hard to collect in real world: depth, normal, optical flow, scene flow, long-term tracking, 6D pose. Standard paradigm: pretrain on synthetic, adapt/fine-tune/evaluate on real. \key{Domain gap}: synthetic accuracy $\ne$ real accuracy; gap comes from rendering, texture, lighting, object distribution, sensor noise, behavior.
\key{Generation paradigms}: 3D/game engines are realistic and physically grounded but asset-limited; procedural generation gives infinite scenes + perfect ground truth but may look synthetic; GenAI looks realistic but geometry/labels may be noisy. Key sources: FlyingThings3D/SceneFlow, Kubric/MOVi, Hypersim, Replica, GTA5/SYNTHIA/AI2-THOR, Virtual KITTI.
\subsection{Tracking \& Pose}
Synthetic tracking provides full trajectories and occlusion labels. Point tracking task: given query point in a frame, predict its location and visibility across time; long-term trajectories and occlusion flags are almost impossible to annotate densely by hand. Kubric uses Blender+PyBullet to produce RGB/depth/normal/segmentation/flow/trajectories/visibility. PIPs: persistent independent particles track any point; TAP-Vid-Kubric benchmark; TAP-Net/TAPIR refine per-frame matching into robust point tracking; CoTracker tracks points jointly with cross-track attention; BootsTAPIR/CoTracker3 use teacher-student pseudo labels; Track-On-R adds verifier-guided filtering.
6-DoF pose: model-based uses CAD model/rendering; model-free/few-shot matches RGB-D/reference observations. FS6D uses few-shot ShapeNet6D, online augmentation (texture blending/deformation). MegaPose: render-and-compare. FoundationPose: unified framework trained purely synthetic. Metrics:
\[
ADD=\frac1{|M|}\sum_{x\in M}\|Rx+t-(R^*x+t^*)\|,
\]
\[
ADD\text{-}S=\frac1{|M|}\sum_{x_1\in M}\min_{x_2\in M}\|Rx_1+t-(R^*x_2+t^*)\|
\]
for symmetric objects.

\subsection{Depth, Normal, Sensors}
Marigold fine-tunes Stable Diffusion for depth; DAC generalizes to arbitrary cameras incl. ERP panorama. Metric depth is absolute meters; relative/affine-invariant depth preserves ordering/shape up to scale/shift. DAC maps arbitrary FoV cameras into an ERP universal canvas via FoV alignment, then maps prediction back. Camera Depth Models (CDM) denoise RGB + raw noisy sensor depth into clean metric depth, not monocular depth. Surface normal: Omnidata uses 3D scans for mid-level cues; normals are ill-posed from images due to shading/texture/lighting ambiguity. DSINE injects geometry bias: ray direction encoding from intrinsics, Ray ReLU visible hemisphere, local normal rotation, iterative refinement. StableNormal uses diffusion prior with YOSO init + SG-DRN refinement. Key lesson: right inductive bias can replace huge data scale.

\section{L11 Curating Vision-Language Data}
\key{Alt-text web data}: large scale but noisy/mismatched. Uses: generation, VLM, VLA. Pipeline: crawl image-text, text filtering, image-text alignment verification, dedup/safety/balancing.
\key{Three schools}: text-only filtering (WIT/MetaCLIP), image-text alignment verification (CC3M/CC12M/LAION/DataComp/DFN), and embracing noise/scale (ALIGN/No Filter/WebLI). Quality controls precision; diversity controls style/long-tail/cultural coverage; safety filtering is still needed even when semantic filtering is relaxed.
\subsection{Text-Side Filtering}
General: build vocabulary/concepts, match text, balance counts. WIT behind CLIP; WIT vs YFCC: scale and web alt-text matter. MetaCLIP demystifies CLIP data: use metadata vocabulary; cap per concept (e.g. 20k) to balance frequent/rare terms; performance comes from transparent curation + scale.
\subsection{Image-Text Alignment}
Filtering paradox: stricter filters remove noise but also diversity/long tail. CC3M: harvest alt-text, image filtering, text filtering, Cloud Vision API tag overlap, hypernymization; high precision but low recall. CC12M keeps image-text verification but relaxes unimodal filters, increasing scale/diversity with lower precision. LAION uses CLIP cosine threshold and releases URLs+metadata; democratizes scale but inherits CLIP filter bias and lacks concept balancing. DataComp: benchmark dataset curation strategies under fixed compute; top-30\% CLIP score is a Western-benchmark sweet spot, not universal optimum; dedup improves generalization. DFN trains a dedicated filter on high-quality image-text pairs; filter quality $\ne$ ImageNet accuracy, and poisoned/low-quality filter training data hurts badly. No Filter/WebLI: at 10B--100B scale, random noise can dilute but systematic bias/stereotypes amplify; global pretraining + short English fine-tune preserves diversity. Scale helps coverage parity and multilingual/cultural performance, not automatic stereotype correction. Datasets still need NSFW/toxicity/PII/watermark/low-quality filtering, but overfiltering harms coverage.

\section{L12 Diffusion \& Score Models}
\subsection{EBM, Score, DSM}
Density modeling wants sample from $p(x)$. EBM: $p_\theta(x)=\exp(-E_\theta(x))/Z_\theta$, $Z$ intractable. Langevin dynamics uses score:
\[
x_{k+1}=x_k+\frac{\eta}{2}\nabla_x\log p(x_k)+\sqrt{\eta}\epsilon_k.
\]
Score matching learns $s_\theta(x)\approx\nabla_x\log p(x)$ without $Z$:
\[
\begin{aligned}
\mathcal J_{\mathrm{SM}}
&=\frac12\E_p\|s_\theta(x)-\nabla_x\log p(x)\|^2\\
&\equiv \E_p\left[\frac12\|s_\theta(x)\|^2+\nabla_x\cdot s_\theta(x)\right]+\text{const}.
\end{aligned}
\]
Hyvarinen score matching removes unknown $\nabla\log p$ by integration by parts; denoising score matching avoids second derivatives by learning scores of noisy conditional Gaussians. With $x_\sigma=x+\sigma\epsilon$:
\[
\nabla_{x_\sigma}\log p_\sigma(x_\sigma|x)=-(x_\sigma-x)/\sigma^2.
\]
NCSN learns scores at multiple $\sigma$; annealed Langevin samples from large noise to small noise. Metropolis-Hastings accepts/rejects random-walk proposals but mixes slowly in high dimensions; Langevin uses the score direction for informed updates. Annealed Langevin: for $\sigma_1>\cdots>\sigma_L$, run several Langevin steps at each noise level, coarse-to-fine.

\subsection{DDPM}
Forward:
\[
q(x_t|x_{t-1})=\N(\sqrt{\alpha_t}x_{t-1},(1-\alpha_t)I),\quad \alpha_t=1-\beta_t,
\]
\[
\begin{gathered}
\bar\alpha_t=\prod_{s=1}^t\alpha_s,\quad
q(x_t|x_0)=\N(\sqrt{\bar\alpha_t}x_0,(1-\bar\alpha_t)I),\\
x_t=\sqrt{\bar\alpha_t}x_0+\sqrt{1-\bar\alpha_t}\epsilon.
\end{gathered}
\]
Reverse:
\[
p_\theta(x_{t-1}|x_t)=\N(\mu_\theta(x_t,t),\Sigma_\theta(x_t,t)).
\]
Posterior:
\[
q(x_{t-1}|x_t,x_0)=\N(\tilde\mu_t,\tilde\beta_tI),
\]
\[
\tilde\mu_t=\frac{\sqrt{\bar\alpha_{t-1}}\beta_t}{1-\bar\alpha_t}x_0+
\frac{\sqrt{\alpha_t}(1-\bar\alpha_{t-1})}{1-\bar\alpha_t}x_t,\quad
\tilde\beta_t=\frac{1-\bar\alpha_{t-1}}{1-\bar\alpha_t}\beta_t.
\]
Noise parameterization:
\[
x_0=\frac{x_t-\sqrt{1-\bar\alpha_t}\epsilon}{\sqrt{\bar\alpha_t}},
\qquad
\mathcal L_{\mathrm{simple}}=\E_{t,x_0,\epsilon}\|\epsilon-\epsilon_\theta(x_t,t)\|_2^2.
\]
Sampling:
\[
x_{t-1}=\frac{1}{\sqrt{\alpha_t}}\left(x_t-\frac{\beta_t}{\sqrt{1-\bar\alpha_t}}\epsilon_\theta(x_t,t)\right)+\sigma_tz.
\]

\subsection{Continuous View \& Equivalent Targets}
SDE/PF-ODE unify score/diffusion:
\[
\begin{gathered}
\dd x=f(x,t)\dd t+g(t)\dd w,\\
\dd x=\left[f(x,t)-\frac12g(t)^2\nabla_x\log p_t(x)\right]\dd t.
\end{gathered}
\]
Gaussian path:
\[
z_t=\alpha_t x+\sigma_t\epsilon,\quad \mathrm{SNR}=\alpha_t^2/\sigma_t^2,\quad \lambda=\log \mathrm{SNR}.
\]
Equivalent targets: clean $x$, noise $\epsilon$, score $s_t(z)=-\E[\epsilon|z_t=z]/\sigma_t$, velocity/path derivative $v_t=\dot\alpha_tx+\dot\sigma_t\epsilon$, denoiser $D(z_t)=\E[x|z_t]$. Weighting and schedule can be expressed in SNR/log-SNR coordinates; epsilon-prediction is not uniform in SNR, log-SNR makes interpretation cleaner.
\key{Forward paths}: VE/EDM $x_t=x_0+\sigma_t\epsilon$; VP/DDPM $x_t=\alpha_tx_0+\sigma_t\epsilon,\alpha_t^2+\sigma_t^2=1$; Flow Matching $x_t=(1-t)x_0+t\epsilon$. Only VP keeps variance normalized. \key{Big equivalence}: diffusion $\approx$ score learning $\approx$ denoising $\approx$ flow matching $\approx$ probability-flow ODE for Gaussian paths. \key{VLB weight}: in $x_0$ prediction, DDPM VLB weights each step by SNR drop,
\[
L_{t-1}\propto (r_{t-1}-r_t)\|D_\theta(x_t,t)-x_0\|^2,\qquad r=\mathrm{SNR}.
\]
\key{Training design}: choose path, prediction target $(\epsilon,x_0,s,v)$, objective weight $w(\lambda)$, and noise proposal/schedule $q(\lambda)$; Monte Carlo estimator uses weight $w/q$. Monotone weighting is not just heuristic: it can be interpreted as a mixture of ELBOs on Gaussian-noise-augmented data. Min-SNR clips high-SNR weights; sigmoid/EDM emphasize mid-noise or detail levels.

\section{L13 Beyond Diffusion}
\subsection{Normalizing Flow}
Flow uses invertible $z=f_\theta(x)$ with exact likelihood:
\[
\log p_X(x)=\log p_Z(f_\theta(x))+\log\left|\det\frac{\partial f_\theta(x)}{\partial x}\right|.
\]
MLE minimizes negative log-likelihood. Composition: log-det sums over layers. Bottleneck: invertibility + tractable Jacobian determinant, dimension match. Triangular Jacobian makes determinant product of diagonals. NICE additive coupling:
\[
y_a=x_a,\quad y_b=x_b+m(x_a),\quad \det J=1.
\]
RealNVP affine coupling:
\[
y_a=x_a,\quad y_b=x_b\odot\exp(s(x_a))+t(x_a),\quad \log|\det J|=\sum s(x_a).
\]
TarFlow: transformer masked autoregressive flow over patches; triangular Jacobian via ordering. Noise augmentation + Tweedie denoising:
\[
\E[x|y]=y+\sigma^2\nabla_y\log p(y).
\]
Guidance analogous classifier-free: combine conditional/unconditional scores/log-probs.

\subsection{GAN, AR, Masked/Flexible-Order}
GAN objective:
\[
\min_G\max_D\E_{x\sim p_{data}}\log D(x)+\E_{z\sim p_z}\log(1-D(G(z))).
\]
Pros: fast high-fidelity samples; cons: unstable, mode collapse, no exact likelihood. Optimal discriminator:
\[
\begin{gathered}
D^*(x)=\frac{p_{data}(x)}{p_{data}(x)+p_G(x)},\\
V(D^*,G)=-\log4+2\,JSD(p_{data}\|p_G).
\end{gathered}
\]
If supports barely overlap, JS can saturate and generator gradients become poor. WGAN replaces JS by Earth-Mover/Wasserstein-1:
\[
W_1(p,q)=\sup_{\|f\|_L\le1}\E_p f(x)-\E_q f(x),
\]
with 1-Lipschitz critic, giving useful gradients even for non-overlapping supports. RpGAN compares paired real/fake scores; R3GAN adds principled R1/R2-style regularization to stabilize dynamics. Modern GANs also survive as perceptual/adversarial losses for diffusion, distillation, and few-step generators. AR factorizes exact likelihood:
\[
p(x)=\prod_{i=1}^N p(x_i|x_{<i}),
\]
stable MLE/cross-entropy but serial $O(N)$ sampling and ordering bias. Image AR tokenizes pixels/patches/codebook tokens; Transformer provides long-range context. \key{AR evolution}: PixelRNN models raw pixels but training/inference are sequential; PixelCNN uses masked convolution so training is parallel over pixels but sampling remains sequential. VQ-VAE shortens image sequence via discrete codebook tokens. Loss has reconstruction + codebook + commitment terms; straight-through estimator copies gradient through nearest-code lookup. JetFormer-style soft tokens use invertible flow tokenizer + AR Transformer, preserving continuous information and exact likelihood more directly than hard codebooks.
Masked/flexible-order generation predicts masked tokens in iterative refinement; faster than AR, useful for editing/inpainting, likelihood often approximate or schedule-dependent. \key{MaskGIT}: VQGAN tokenizer + bidirectional Transformer; train by masking random token subset and cross-entropy on masked tokens only. Iterative decoding predicts all masked positions in parallel, samples tokens, keeps high-confidence ones, and follows a mask schedule; steps drop from $O(N)$ to about $8$--$12$, but bidirectional attention cannot use KV cache. \key{RandAR}: random-order causal generation adds position instruction tokens, allowing KV-cache and parallel decoding of multiple requested positions while avoiding fixed raster order. \key{Comparison}: Flow exact likelihood/fast/stable but Jacobian constraints; GAN no likelihood/fast but unstable and mode collapse; AR exact likelihood/stable but slow serial order; Masked approx likelihood/medium speed with mask schedule; Diffusion ELBO/score/stable but many denoise steps.

\section{Cross-Lecture Exam Hooks}
\key{Geometry chain}: pinhole projection $\to$ depth ambiguity $\to$ stereo disparity $\to$ epipolar constraint $\to$ rectification $\to$ dense depth. \key{Feature chain}: local keypoints/SIFT/RANSAC $\to$ homography/mosaic/SfM; modern features from SSL/VLM. \key{CNN-to-Transformer chain}: conv local bias and weight sharing; attention global interaction; ViT patch tokens; ConvNeXt borrows Transformer design. \key{Dense chain}: FCN $\to$ U-Net/FPN $\to$ Mask R-CNN/SAM/depth/diffusion. \key{Data chain}: supervised labels expensive; SSL, synthetic data, web data curation, pseudo-label teacher/student all scale supervision.

\end{multicols*}
\end{document}
